\chapter{Market Making: Quoting, Inventory, and Avellaneda--Stoikov}
\label{ch:market_making}
\chaptermark{Market making and Avellaneda--Stoikov}

You are quoting RAINFOREST\_RESIN on the Prosperity exchange. The
product's fair value sits around $10\,000$ seashells, barely moves,
and every other team on the leaderboard is printing steady profit off
it. You can post any limit bid at any tick and any limit ask at any
tick. What do you quote?

A first guess: bid one tick below $10\,000$, ask one tick above,
collect the spread forever. The guess is almost right (top
Prosperity teams quote something close to it), but it has to survive
two adversaries already met in \cref{ch:participants,ch:kyle_gm}: a
counterparty who knows more than you, and flow that drifts the price
against whatever inventory you hold. The bid-ask spread a maker earns
is precisely the compensation demanded for bearing those two
risks: \emph{adverse selection} (trading against better-informed
counterparties) and \emph{inventory risk} (carrying unwanted positions
through a volatile mid). A naive symmetric quoter beats neither. This chapter builds from first principles the machinery that
fixes it: the \textbf{Avellaneda--Stoikov (A--S) market-making model}
\citeAS{}, the closed-form rule for how a risk-averse quoter skews
bid and ask as inventory drifts. The route is: diagnose the naive
quote against the \cref{ch:lob} toy book; derive the
\textbf{reservation price} $r = \Spx - \invq\riskav\volat^2(T-t)$
from an indifference argument (stating the Hamilton--Jacobi--Bellman
equation but not solving it); state the \textbf{optimal spread}
$\delta^* = \riskav\volat^2(T-t) + (2/\riskav)\ln(1 + \riskav/\arrate)$
with its two-term decomposition and sanity limits; and close with the
failure modes, where A--S does \emph{not} apply.

\begin{backgroundbox}[Prerequisites]
From earlier in the book:
\begin{itemize}
  \item \Cref{ch:lob}: the limit order book, limit versus market
  orders, the maker/taker distinction, the matching engine, and
  the specific four-level toy LOB with best bid
  $\bidp = 102.54$ (quantity $131$), best ask
  $\askp = 103.23$ (quantity $48$), next bid $101.87$ (quantity
  $32$), next ask $103.98$ (quantity $84$), and mid
  $\midp = 102.885$ with spread $\spreads = 0.69$.
  \item \Cref{ch:participants}: the trader taxonomy, the
  informed/uninformed split, and Olivia the informed bot with
  $\alpha = 0.3$ and $V \in \{100, 110\}$.
  \item \Cref{ch:price_value}: the formal definitions of bid,
  ask, mid, microprice, and the distinction between a quoted
  mid and a true fair value.
  \item \Cref{ch:spreads_roll}: the three-cause decomposition of
  the spread (order-processing, inventory, adverse-selection).
  Each of these three causes will reappear as a parameter in
  A--S, which is this chapter's unifying theme.
  \item \Cref{ch:kyle_gm}: Kyle's linear equilibrium with
  price impact $\lamk$ and the Glosten--Milgrom equilibrium
  spread as an adverse-selection premium, the formal
  machinery behind why a zero-spread market maker loses money.
  \item \Cref{ch:market_impact}: the empirical square-root
  impact law and the distinction between permanent and
  temporary impact.
  \item \Cref{ch:order_flow}: order-flow imbalance, queue
  position as a source of alpha, the qualitative Hawkes
  clustered-arrivals picture, and the economic view of latency
  as a second-derivative game. A--S's arrival intensity
  $\arrate$ is calibrated directly from order-flow data of the
  kind that chapter introduced.
\end{itemize}
From outside the book:
\begin{itemize}
  \item Elementary probability: expectation, variance, Poisson
  processes with rate $\lambda$, the memoryless property.
  \item Elementary calculus: second-order Taylor expansion, and
  basic single-variable differentiation for optimisation. No
  stochastic calculus needed beyond the statement of the value
  function; the HJB equation is stated, not solved.
  \item Basic utility theory: exponential (CARA) utility
  $U(X) = -\exp(-\riskav X)$, with $\riskav > 0$ the risk
  aversion. The reader only needs to know it is a concave
  function of wealth.
\end{itemize}
\end{backgroundbox}

\section{The naive symmetric market maker}
\label{sec:ch8_naive}

Start with the simplest policy: pick a half-width $\delta$, quote a
bid at $\midp - \delta$ and an ask at $\midp + \delta$, wait. Each
arriving buyer lifting the ask earns $\delta$ above the mid; each
seller hitting the bid earns $\delta$ below. The round-trip gain per
matched buy--sell pair is $2\delta$, the \emph{spread capture}.

In a world of symmetric uninformed flow (equal Poisson buyers and
sellers, no mid drift, no memory), this policy is optimal and
profits grow linearly. \Cref{ch:participants} called that world
fiction. The real world has two structural deviations, either of
which can dominate spread capture on any single run; the numerical
example below, on the toy LOB of \cref{ch:lob}, makes both concrete.

\subsection{What the naive maker earns: a benign run}
\label{subsec:naive_benign}

Take the \cref{ch:lob} LOB. The best bid is $\bidp = 102.54$
(quantity $131$), the best ask is $\askp = 103.23$ (quantity $48$),
and the mid is $\midp = 102.885$. The round-trip spread is
$\spreads = 0.69$, giving a half-spread of $\delta = 0.345$. Set
the market maker's quotes to coincide with the top of book:
$\bidp_{\mathrm{mm}} = 102.54$ and $\askp_{\mathrm{mm}} = 103.23$.
(The round numbers are deliberately chosen so the toy sits exactly
at the inside of the book; for this illustration we assume the
market maker always wins queue priority on their side.)

Ten customer orders arrive. Six are buys (each lifts one share of
the market maker's ask); four are sells (each hits one share of the
market maker's bid). After the ten arrivals, the market maker has:
\begin{align*}
  \text{sold at ask: } & 6 \text{ shares} \times 103.23
      &&= +\$619.38 \text{ in cash} \\
  \text{bought at bid: } & 4 \text{ shares} \times 102.54
      &&= -\$410.16 \text{ in cash} \\
  \text{net cash: } & && +\$209.22 \\
  \text{net inventory: } & \underbrace{-6}_{\text{sold}} +
      \underbrace{+4}_{\text{bought}}
      &&= -2 \text{ shares (net short)}
\end{align*}

The market maker ends the batch \emph{short two shares}. To compare
this to the spread capture alone, subtract off the mark-to-market
value of the residual inventory at the mid:
\begin{align}
  \text{edge earned}
  \;&=\; \text{cash} \;-\; \invq \cdot \midp \\
  \;&=\; 209.22 \;-\; (-2)(102.885) \\
  \;&=\; 209.22 + 205.77 \;=\; 414.99. \label{eq:ch8_edge_mismatch}
\end{align}
That is wealth in absolute cash, which does not match the intuitive
``earned the spread'' number. The cleaner decomposition measures
earnings \emph{per share crossed}: every customer crossing earns the
maker $\midp - \bidp_{\mathrm{mm}} = 0.345$ or $\askp_{\mathrm{mm}}
- \midp = 0.345$ depending on direction. Across the ten arrivals,
\begin{equation}
  \text{edge} \;=\; 6 \cdot (\askp_{\mathrm{mm}} - \midp)
           \;+\; 4 \cdot (\midp - \bidp_{\mathrm{mm}})
  \;=\; 10 \cdot 0.345 \;=\; 3.45.
  \label{eq:ch8_edge}
\end{equation}
The maker has earned $3.45$ seashells of edge and is carrying a
residual short of $-2$ shares. If the mid drifts up by $0.5$, the
inventory marks down by $-2 \cdot (+0.5) = -1.0$ seashells. Total
batch P\&L:
\begin{equation}
  \underbrace{3.45}_{\text{edge}}
  \;+\; \underbrace{(-2)(0.5)}_{\text{inventory drift}}
  \;=\; 2.45 \text{ seashells}.
  \label{eq:ch8_benign_net}
\end{equation}
The market maker keeps most of what they earned. On any single
benign run the model looks fine.

\subsection{What the naive maker loses: a directional run}
\label{subsec:naive_bad}

Rerun the same policy on less friendly flow. Ten orders arrive, all
buys: every one lifts the ask, none hits the bid. The maker sells
ten shares into the buying crowd and ends ten shares \emph{short}.
Meanwhile, and this is the key point, the mid drifts up by $2$
seashells as buying pressure carries the market with it. The decomposition:
\begin{align}
  \text{edge}
  \;&=\; 10 \cdot (\askp_{\mathrm{mm}} - \midp)
  \;=\; 10 \cdot 0.345 \;=\; 3.45, \\[0.25em]
  \text{inventory drift}
  \;&=\; (-10)(+2.0) \;=\; -20.0, \\[0.25em]
  \text{net P\&L}
  \;&=\; 3.45 \;-\; 20.0 \;=\; -16.55. \label{eq:ch8_bad_net}
\end{align}
The maker collected $3.45$ seashells of edge and was handed back
$-20$ in mark-to-market loss on the unhedged short: $-16.55$
underwater \emph{and} still carrying the exposure. The arithmetic
is symmetric under a sign flip: ten sells and a downward drift would
leave the maker net long and equally underwater.

\begin{pitfallbox}[Symmetric quoting kills you on trend]
The naive symmetric market maker's P\&L has two terms. The
\emph{edge} term scales linearly in the number of fills and in
the half-spread $\delta$: it grows at most like
$O(\text{fills} \cdot \delta)$. The \emph{inventory-drift} term
scales like $\invq \cdot \Delta \midp$, which can grow as fast as
(number of fills) $\times$ (total mid move) when flow is
one-sided. On any persistent directional run (trend,
newsbreak, stop cascade, or informed flow in the
\cref{ch:kyle_gm} sense), the second term wins and the maker
loses more on the inventory than they collected on the spread.
This is the entire economic content of the observation
``market makers lose to trends.'' The fix is not to widen
$\delta$ (which reduces fill rate linearly); it is to
\emph{skew} the quotes against inventory so the fills
themselves mean-revert the position.
\end{pitfallbox}

The run also exposes adverse selection in the sense of
\cref{ch:kyle_gm}: the ten buyers were not drawn from a zero-drift
distribution; they knew something the maker did not. The
Glosten--Milgrom calculation of \cref{sec:gm} predicts a symmetric
quoter loses at least $\alpha(V_H - V_L)/2$ per fill against informed
counterparties; on this run every counterparty was informed. The
loss shows up as the $-20$ seashells of drift, Kyle's $\lamk$ acting
on the flow.

A maker who refuses to skew against inventory is exposed to both
failures. A--S fixes both through a single knob, the risk-aversion
coefficient $\riskav$.

\section{Inventory-aware quoting: the reservation price}
\label{sec:ch8_reservation}

The previous section's lesson: quotes should not be centred on the
market mid, but on a price that \emph{depends on current inventory}.
Long inventory: shift both quotes down so the next matching flow is
more likely a buy (selling the long down). Short inventory: shift
both up so the next flow is more likely a sell (buying the short
back). In A--S language, quotes are centred on a private,
inventory-dependent \textbf{reservation price}.

\subsection{Setup}
\label{subsec:ch8_setup}

A single risky asset has mid $\Spx_t$ evolving as driftless
Brownian motion
\begin{equation}
  \mathrm{d}\Spx_t \;=\; \volat \, \mathrm{d}W_t,
  \label{eq:ch8_sde}
\end{equation}
with constant $\volat > 0$ and $W_t$ a standard Brownian motion.
The maker quotes bid $\bidp_t$ and ask $\askp_t$, holds inventory
$\invq_t \in \Z$, and has cash $X_t$. Fills arrive as Poisson
processes with intensity specified in \cref{sec:ch8_spread}; the
reservation-price derivation does not yet need the arrival rate.
Trading horizon $[0, \Tmat]$, with terminal inventory marked to the
mid:
\begin{equation}
  \text{terminal wealth}
  \;=\; X_\Tmat \;+\; \invq_\Tmat \cdot \Spx_\Tmat.
\end{equation}

The market maker is risk-averse with exponential (CARA) utility
\begin{equation}
  U(W) \;=\; -\exp(-\riskav W),
  \qquad \riskav > 0.
\end{equation}
CARA utility has the convenient property that optimal behaviour
depends only on risk aversion $\riskav$, not on initial wealth.
Define the \textbf{value function}
\begin{equation}
  v(t, x, \Spx, \invq)
  \;\defeq\; \sup_{\bidp, \askp}
    \E\!\left[ -\exp\!\bigl(-\riskav (X_\Tmat + \invq_\Tmat \Spx_\Tmat)\bigr)
       \;\Big|\; X_t = x,\, \Spx_t = \Spx,\, \invq_t = \invq \right],
  \label{eq:ch8_valuefn}
\end{equation}
with supremum over admissible quoting policies. The value function
encodes ``best expected utility achievable from state
$(x, \Spx, \invq)$ at time $t$.''

\subsection{The HJB equation, stated}
\label{subsec:ch8_hjb}

Stochastic optimal control turns \eqref{eq:ch8_valuefn} into a PDE.
Formally \citeAS{, §3}, \citeCartea{, Ch 10.2}, $v$ satisfies the
Hamilton--Jacobi--Bellman equation
\begin{align}
  \partial_t v
  \;+\;\tfrac{1}{2}\volat^2 \, \partial_\Spx^2 v
  \;&+\; \max_{\bidp}
      \lambda^b(\Spx - \bidp)\bigl[\,v(t, x - \bidp, \Spx, \invq + 1) - v\,\bigr]
  \nonumber \\
  \;&+\; \max_{\askp}
      \lambda^a(\askp - \Spx)\bigl[\,v(t, x + \askp, \Spx, \invq - 1) - v\,\bigr]
  \;=\; 0,
  \label{eq:ch8_hjb}
\end{align}
with terminal condition
\begin{equation}
  v(\Tmat, x, \Spx, \invq)
  \;=\; -\exp\!\bigl(-\riskav (x + \invq \Spx)\bigr).
  \label{eq:ch8_terminal}
\end{equation}
The first line is mid diffusion; the next two are the control
terms, one per side, each reading ``pick the quote maximising
expected utility gain from a fill, weighted by Poisson fill
intensity.'' The intensities $\lambda^b, \lambda^a$ depend on
distance from the mid and are parameterised in
\cref{sec:ch8_spread}.

This chapter will \textbf{not} solve \eqref{eq:ch8_hjb}.
\citet{avellaneda2008high} solve it via a substitution separating
the $(x, \Spx)$ dependence from the inventory dependence, yielding a
one-dimensional ODE in $\invq$ with an exact solution. The three
pages of algebra add nothing to intuition; only the output matters
here.

\begin{remark}[Why state an unsolved HJB?]
The reservation-price and optimal-spread formulas are genuine
byproducts of a stochastic-control problem, not identities one can
guess. Stating the HJB shows where the closed form comes from and
points readers at \citeAS{} and \citeCartea{} for the algebra. The
same pattern recurs in \cref{ch:almgren_chriss} for optimal
execution.
\end{remark}

\subsection{Reservation price via indifference}
\label{subsec:ch8_indifference}

A shorter route avoids the HJB entirely and is the one worth
internalising. Define the reservation price by an
\textbf{indifference} condition: $r(t, \Spx, \invq)$ is the price at
which a maker holding $\invq$ shares is indifferent between the
current portfolio and the portfolio after buying one more share at
$r$. Equivalently, it is the private mid the maker would use if
forced to quote a single number and trade both sides at it.

Write the indifference condition in utility-adjusted wealth. With
cash $x$ and inventory $\invq$, the certainty-equivalent portfolio
value is
\begin{equation}
  \mathrm{CE}(t, x, \Spx, \invq)
  \;=\; -\frac{1}{\riskav} \ln\!\bigl(-\E[\,U(X_\Tmat +
         \invq_\Tmat \Spx_\Tmat)\,]\bigr),
\end{equation}
with expectation under the pure-hold strategy: no trading between
$t$ and $\Tmat$, just mark $\invq$ to the terminal mid. Under
\eqref{eq:ch8_sde} the terminal mid is
$\Spx_\Tmat \sim \Normal(\Spx, \volat^2(\Tmat - t))$, so terminal
wealth is normal with mean $x + \invq \Spx$ and variance
$\invq^2 \volat^2 (\Tmat - t)$. For CARA utility and a normal
argument, the standard identity
\begin{equation}
  \E[-\exp(-\riskav Y)]
  \;=\; -\exp\!\bigl(-\riskav\,\E[Y] + \tfrac{1}{2}\riskav^2 \Var[Y]\bigr)
  \label{eq:ch8_cara_normal}
\end{equation}
applies, so the certainty equivalent collapses to
\begin{equation}
  \mathrm{CE}(t, x, \Spx, \invq)
  \;=\; x \;+\; \invq \Spx
      \;-\; \tfrac{1}{2}\riskav \invq^2 \volat^2 (\Tmat - t).
  \label{eq:ch8_ce}
\end{equation}
The first two terms are the mark-to-market value; the third is the
\emph{risk penalty} the inventory-$\invq$ holder pays for
mark-to-market variance between now and $\Tmat$. It is
negative-definite in $\invq$: bigger positions hurt more.

The reservation price is now defined by the indifference equation:
the maker is indifferent between $(x, \invq)$ and $(x - r, \invq + 1)$,
i.e.\ paying $r$ to acquire one more share:
\begin{equation}
  \mathrm{CE}(t, x, \Spx, \invq)
  \;=\; \mathrm{CE}(t, x - r, \Spx, \invq + 1).
  \label{eq:ch8_indiff}
\end{equation}
Plug \eqref{eq:ch8_ce} into both sides. All linear-in-$\Spx$ and
linear-in-$x$ terms cancel except the payment $r$ and the shift of
the quadratic penalty. After cancellation,
\begin{align}
  \invq \Spx - \tfrac{1}{2}\riskav \invq^2 \volat^2 (\Tmat-t)
  \;&=\;
  -r + (\invq+1)\Spx
  - \tfrac{1}{2}\riskav (\invq+1)^2 \volat^2(\Tmat-t),
\end{align}
\begin{align}
  r \;&=\; \Spx \;-\; \tfrac{1}{2}\riskav \volat^2 (\Tmat-t)\,
          \bigl[(\invq+1)^2 - \invq^2\bigr] \nonumber\\
  \;&=\; \Spx \;-\; \tfrac{1}{2}\riskav \volat^2 (\Tmat-t)(2\invq+1).
\end{align}
For the marginal (infinitesimal) indifference defining the A--S
continuous-share reservation price, drop the $+1$ relative to
$2\invq$ (equivalent to differentiating the certainty equivalent
in $\invq$):
\begin{equation}
  \boxed{\;
    r(t, \Spx, \invq)
    \;=\; \Spx \;-\; \invq \, \riskav \, \volat^2 \,(\Tmat - t)\;}
  \label{eq:ch8_r_closedform}
\end{equation}
which is the Avellaneda--Stoikov reservation price
\citeAS{, eq.\ (9)}. The full HJB solution reproduces this formula
with no correction terms in the large-$\arrate\volat^2\Tmat/\riskav$
regime that \citeCartea{} flag as practically-relevant. The
indifference derivation above is not an approximation but the
exact marginal statement.

\begin{keyfactbox}[Avellaneda--Stoikov reservation price]
The A--S reservation price is
\[
  r(t, \Spx, \invq) \;=\; \Spx \;-\; \invq \, \riskav \, \volat^2 \,(\Tmat - t).
\]
Reading the parameters:
\begin{itemize}
  \item $\Spx$ is the current market mid.
  \item $\invq$ is the market maker's current inventory in shares
  (positive $=$ long, negative $=$ short).
  \item $\riskav > 0$ is the market maker's risk aversion.
  \item $\volat$ is the mid's volatility (same units as $\Spx$,
  per $\sqrt{\text{time}}$).
  \item $\Tmat - t$ is the time remaining until the risk horizon
  (end of session, end of trading day).
\end{itemize}
The second term is the \emph{inventory skew}: the shift of the
reservation price away from the market mid in response to an
open position. Long inventory ($\invq > 0$) pushes $r$ \emph{down}
relative to $\Spx$; short inventory ($\invq < 0$) pushes $r$
\emph{up}. The magnitude of the shift grows with risk aversion,
with variance per unit time $\volat^2$, and with time remaining
$(\Tmat - t)$.
\end{keyfactbox}

\begin{intuitionbox}
The reservation price is your \emph{private} fair value. It is
what you would accept to trade one more share in either
direction, given that you already hold $\invq$ shares and are
averse to carrying position risk over the remaining horizon. A
risk-neutral trader ($\riskav = 0$) would have $r = \Spx$: they
are equally happy to hold five shares or fifty. A risk-averse
trader with non-zero inventory knows the variance of their
mark-to-market is $\invq^2 \volat^2 (\Tmat - t)$, and shifts
their private mid in the direction that makes the next trade
reduce $\invq^2$. A--S centres the market maker's \emph{quotes}
on this private mid, not on the market mid.
\end{intuitionbox}

\begin{figure}[ht]
\centering
\begin{tikzpicture}[scale=1.1]
\draw[thick, -Stealth] (-5.2,0) -- (5.8,0) node[right, font=\small]{Price};
% Key points
\fill[blue!70]  (-2.4,0) circle (3pt);
\fill[black!50] (-0.6,0) circle (3pt);
\fill[gray!50]  ( 0.0,0) circle (2pt);
\fill[red!70]   ( 2.4,0) circle (3pt);
% Labels above
\node[above=6pt, font=\small, blue!80]  at (-2.4,0) {Bid $r-\delta^b$};
\node[above=6pt, font=\small, black!60] at (-0.6,0)
    {\shortstack{Reservation\\price $r$}};
\node[above=6pt, font=\small, gray]     at ( 0.0,0) {Mid $s$};
\node[above=6pt, font=\small, red!80]   at ( 2.4,0) {Ask $r+\delta^a$};
% Braces below
\draw[decorate, decoration={brace, mirror, amplitude=6pt}]
    (-2.4,-0.15) -- (0,-0.15)
    node[midway, below=8pt, font=\footnotesize]{$\delta^b$};
\draw[decorate, decoration={brace, mirror, amplitude=6pt}]
    (0,-0.15) -- (2.4,-0.15)
    node[midway, below=8pt, font=\footnotesize]{$\delta^a$};
\draw[decorate, decoration={brace, amplitude=6pt}]
    (-2.4,0.15) -- (2.4,0.15)
    node[midway, above=12pt, font=\footnotesize]{spread $\delta^a+\delta^b$};
\end{tikzpicture}
\caption{Avellaneda--Stoikov quoting structure. The maker posts bid at
$r - \delta^b$ and ask at $r + \delta^a$, where $r$ is the
\emph{reservation price}, the mid adjusted downward (upward) when
inventory is long (short). When the maker holds excess long inventory,
$r$ shifts below the market mid, making the bid less competitive and
the ask more competitive, steering inventory back to zero.}
\label{fig:as_number_line}
\end{figure}

\subsection{A worked reservation-price example, and a calibration lesson}
\label{subsec:ch8_r_worked}

Take a concrete Prosperity-style setup. The product is
RAINFOREST\_RESIN, the market mid is $\Spx = 10\,000$ seashells, the
market maker currently holds $\invq = +10$ shares (a long position),
volatility is $\volat = 5$ seashells per unit time, and
$\Tmat - t = 100$ units. A careless first choice of the risk
aversion parameter is $\riskav = 0.1$. Plug in:
\begin{align}
  r \;&=\; 10\,000 \;-\; (+10)\cdot 0.1 \cdot 25 \cdot 100
  \nonumber\\
  \;&=\; 10\,000 \;-\; 2500
  \;=\; 7500.
  \label{eq:ch8_r_bad}
\end{align}
The formula says a maker holding ten shares long should shift the
private mid to $7500$, a $2500$-seashell shift. That is absurd:
no maker would lower quotes by that much to offload ten shares on a
product whose round-trip spread is a few seashells.

The absurdity is a calibration lesson, not a bug. The dimensioned
quantity controlling inventory skew is
\begin{equation}
  \riskav \volat^2 (\Tmat - t),
\end{equation}
which has units of (price per share). A--S parameterises the
penalty through $\riskav$; $\volat$ and $\Tmat - t$ are given by the
problem ($\volat$ measured from data, horizon fixed by session). The
only dial the practitioner turns is $\riskav$, and it must be tuned
so that $\riskav \volat^2 (\Tmat - t)$ has the same order of
magnitude as the spread, not orders of magnitude larger.

Rerun the same example with $\riskav = 0.001$:
\begin{align}
  r \;&=\; 10\,000 \;-\; 10 \cdot 0.001 \cdot 25 \cdot 100
  \nonumber\\
  \;&=\; 10\,000 \;-\; 25
  \;=\; 9975.
  \label{eq:ch8_r_good}
\end{align}
The reservation price shifts by $25$ seashells, enough to lean
noticeably against the long inventory but not so much that the
quotes walk off one side of the book.

\begin{pitfallbox}[The product $\riskav \volat^2 (\Tmat - t)$ is what matters]
$\riskav$ is dimensioned (1/currency); only the product
$\riskav \volat^2 (\Tmat - t)$ is economically meaningful.
$\volat$ and $\Tmat - t$ are given by the market and session, so
$\riskav$ is the only real dial. Rule of thumb
\citep{hummingbotASguide}: tune $\riskav$ so that
$\riskav \volat^2 (\Tmat - t)$ is comparable to the typical
per-share spread capture of the symmetric policy at $\invq = 0$.
That keeps the inventory skew large enough to control position
drift but small enough that quotes stay on the book. Copying
$\riskav = 0.1$ from another asset's config and applying it on a
flat-valued product like RAINFOREST\_RESIN is the fastest way to get
a bot that never trades.
\end{pitfallbox}

\section{The optimal bid and ask}
\label{sec:ch8_spread}

The reservation price says \emph{where} to centre quotes, not
\emph{how wide}. Quotes near $r$ fill on almost every arrival and
collect little per fill; wide quotes fill rarely but collect more
per fill. The right half-width maximises the product of fill
probability and edge per fill.

A--S model the fill probability through a Poisson arrival whose
intensity depends on distance from the \emph{market mid}.
Empirically \citeAS{, §4}, the fraction of arriving market orders
reaching a quote at distance $\delta$ decays exponentially. Writing
$\delta^b, \delta^a > 0$ for the distances from mid to bid and ask,
\begin{equation}
  \lambda^b(\delta^b) \;=\; A \exp(-\arrate \, \delta^b),
  \qquad
  \lambda^a(\delta^a) \;=\; A \exp(-\arrate \, \delta^a),
  \label{eq:ch8_arrival}
\end{equation}
where $A > 0$ is the overall market-order arrival rate and
$\arrate > 0$ is the \textbf{decay constant} of fill probability
with distance. Larger $\arrate$ corresponds to a deeper, more liquid
book: moving the quote outward is penalised harder.

\begin{remark}[Origin of the exponential fit]
The $\exp(-\arrate \delta)$ fit is not ad-hoc. A--S derive it from
(i) power-law order sizes (consistent with the impact data of
\cref{ch:market_impact}) and (ii) approximately constant LOB density
around the mid. Under both, the fraction of arriving market orders
reaching a quote at distance $\delta$ decays exponentially, with
$\arrate$ = (power-law exponent) / (mean order size in ticks).
\end{remark}

\subsection{Setup and the optimisation problem}
\label{subsec:ch8_optim_setup}

Extend the \cref{subsec:ch8_setup} setup with fills governed by
\eqref{eq:ch8_arrival}. A bid fill jumps $(x, \invq) \to
(x - \bidp, \invq + 1)$; an ask fill jumps $(x, \invq) \to
(x + \askp, \invq - 1)$. Between fills, the mid diffuses. The
maximisation over $\bidp, \askp$ inside the HJB
\eqref{eq:ch8_hjb} defines the optimal spread.

\citet{avellaneda2008high} solve the HJB via the substitution
$v(t, x, \Spx, \invq) = -\exp(-\riskav(x + \invq\Spx))\,
\theta(t, \invq)^{-\riskav/\arrate}$, turning the PDE into a linear
ODE in $\theta$; \citet{cartea2015algorithmic} give the same
derivation in cleaner notation. The output is closed-form optimal
quotes and, crucially for teaching, an asymptotic formula for
the \emph{total} bid-ask spread $\delta^* = \delta^a + \delta^b$
holding in the regime $\arrate \volat^2 \Tmat / \riskav \gg 1$.

\begin{keyfactbox}[Avellaneda--Stoikov optimal spread]
Under the setup of \cref{subsec:ch8_setup}, with
Brownian mid, exponential fill intensities
$\lambda^{a,b}(\delta) = A e^{-\arrate \delta}$, and CARA utility
with risk aversion $\riskav$, the optimal total quoted spread
(bid to ask) is \citeAS{, eq.\ (17)}, \citeCartea{, §10.4}
\[
  \delta^{*} \;=\; \riskav\,\volat^2\,(\Tmat - t)
             \;+\; \frac{2}{\riskav}\,
               \ln\!\left(1 + \frac{\riskav}{\arrate}\right).
\]
The optimal quotes are placed symmetrically around the
reservation price of \eqref{eq:ch8_r_closedform}:
\[
  p^{*,\,\mathrm{bid}} \;=\; r - \tfrac{1}{2}\delta^{*},
  \qquad
  p^{*,\,\mathrm{ask}} \;=\; r + \tfrac{1}{2}\delta^{*}.
\]
Note the convention: $\delta^*$ is the \emph{total} bid-ask
spread. The half-spread used on each side is $\delta^*/2$.
\end{keyfactbox}

\subsection{Reading the two terms}
\label{subsec:ch8_delta_terms}

The formula has two separable components.

\paragraph{Inventory-risk term, $\riskav \volat^2 (\Tmat - t)$.}
Independent of $\arrate$; widens with risk aversion, volatility, and
time remaining. Economic meaning: the more variance per unit time
and the more time remaining, the more the maker demands for holding
any position across the uncertain window. At $\riskav = 0$
(risk-neutral) the term vanishes. This is the descendant of the
``inventory'' cause of the \cref{ch:spreads_roll} decomposition.

\paragraph{Fill-rate term, $(2/\riskav)\ln(1 + \riskav/\arrate)$.}
Independent of $\volat$ and the horizon; widens as $\arrate$
\emph{decreases} (thinner book). As $\arrate \to \infty$ (infinitely
liquid), the term vanishes: the maker earns the same expected
edge at $\delta = 0$. As $\arrate \to 0$, the term grows without
bound. This is the descendant of the ``order-processing / fill-rate''
cause: to get filled at all, the maker must quote at positive
distance from the mid. The Roll and GM decompositions in
\cref{ch:spreads_roll,ch:kyle_gm} implicitly contain this effect in
their reduced-form ``spread'' parameter.

\subsubsection{Sanity check: the risk-neutral limit}

Check the fill-rate term as $\riskav \to 0$. Expand
$\ln(1 + \riskav/\arrate)$ in powers of $\riskav/\arrate$:
\begin{equation}
  \ln\!\left(1 + \frac{\riskav}{\arrate}\right)
  \;=\; \frac{\riskav}{\arrate}
        \;-\; \tfrac{1}{2}\left(\frac{\riskav}{\arrate}\right)^{\!2}
        \;+\; O\!\left(\frac{\riskav^3}{\arrate^3}\right).
\end{equation}
Multiply by $2/\riskav$:
\begin{equation}
  \frac{2}{\riskav}\,\ln\!\left(1 + \frac{\riskav}{\arrate}\right)
  \;=\; \frac{2}{\arrate}
        \;-\; \frac{\riskav}{\arrate^2}
        \;+\; O(\riskav^2).
\end{equation}
The $\riskav \to 0$ limit gives $2/\arrate$; the inventory-risk
term vanishes too. So a risk-neutral maker quotes
\begin{equation}
  \delta^*_{\mathrm{risk-neutral}} \;=\; \frac{2}{\arrate}.
\end{equation}
Not a coincidence: in the risk-neutral benchmark the maker
independently optimises each side by maximising
$\delta \cdot A\exp(-\arrate\delta)$ in $\delta$, whose unique
maximum is at $\delta = 1/\arrate$; two sides give $2/\arrate$
total. A--S correctly reduces to the independent-sides
profit-maximiser when risk aversion is turned off.

\begin{intuitionbox}
The A--S total spread is the sum of two terms you would write down
anyway. The fill-rate term is how wide you must quote to extract
profit from a one-sided auction ($2/\arrate$, tweaked by $\riskav$).
The inventory term is the extra width for carrying inventory across
a volatile horizon. At $\riskav = 0$ the inventory term vanishes; at
$\volat = 0$ or $\Tmat = t$ it vanishes too. The real work the model
does is the \emph{ratio} between the two terms: it tells you which
failure mode (narrow spreads eaten by variance, or wide spreads
going unfilled) dominates, so you size them against each other.
\end{intuitionbox}

\subsection{Connecting the two terms to the three causes of the spread}
\label{subsec:ch8_three_causes}

\Cref{ch:spreads_roll} decomposed the observed spread into three
causes: \emph{order processing}, \emph{inventory}, \emph{adverse
selection}. A--S maps cleanly onto two of them; the third is ceded
to other models.

\paragraph{Inventory cause.} The A--S inventory-risk term
$\riskav\volat^2(\Tmat - t)$ is the descendant of Roll's inventory
cause. Roll only observed that the spread must contain an
inventory-risk component; A--S pins down the size, proportional to
variance per unit time and risk aversion. \citeCartea{} call this
the ``internal hedging cost'': what the maker would pay to
instantaneously offload one unit of inventory.

\paragraph{Order-processing cause (fill-rate term).} The A--S
fill-rate term $(2/\riskav)\ln(1 + \riskav/\arrate)$ descends from
the order-processing cause. The modern reading is broad: order
processing is the cost of being in the game at all. Quote at the
mid and you fill too often to monetise edge; quote far and you do
not fill. The single-sided optimum is $1/\arrate$; two sides give
$2/\arrate$ as the risk-neutral order-processing contribution.

\paragraph{Adverse-selection cause.} Not in the A--S formula. The
symmetric fill intensities \eqref{eq:ch8_arrival} rule it out by
construction. A--S concedes the adverse-selection component to
Glosten--Milgrom \citeGM{} and Kyle \citeKyle{} (see
\cref{ch:kyle_gm}); the practical fix is to shift the reservation
price with an OFI-based drift estimator (see
\cref{sec:ch8_breaks}). \citeBouchaud{} gives the combined
inventory-plus-adverse-selection quoting rule at length.

Together, the three-cause decomposition and A--S answer ``what is
the spread for?'': the maker must recover order processing
($2/\arrate$), inventory risk ($\riskav\volat^2(\Tmat-t)$), and
adverse selection (ceded to GM/Kyle). Each component has a separate
analytic form, so a well-engineered system treats them as three
modules combined at quoting time.

\subsection{A numerical example, step by step}
\label{subsec:ch8_delta_numeric}

Take the same Prosperity-style setup as in \cref{subsec:ch8_r_worked}:
$\Spx = 10\,000$, $\invq = +10$, $\volat = 5$, $\Tmat - t = 100$,
and now $\arrate = 1.5$. Use the calibrated $\riskav = 0.001$ so
the reservation price comes out to the plausible $r = 9975$
computed in \eqref{eq:ch8_r_good}.

\paragraph{Step 1: inventory-risk term.}
\begin{equation}
  \riskav \volat^2 (\Tmat - t)
  \;=\; 0.001 \cdot 5^2 \cdot 100
  \;=\; 0.001 \cdot 2500
  \;=\; 2.5.
\end{equation}

\paragraph{Step 2: fill-rate term.}
\begin{align}
  \frac{2}{\riskav}\,
    \ln\!\left(1 + \frac{\riskav}{\arrate}\right)
  \;&=\; \frac{2}{0.001}\,
         \ln\!\left(1 + \frac{0.001}{1.5}\right) \nonumber\\
  \;&=\; 2000 \cdot \ln(1.000\overline{6}) \nonumber\\
  \;&=\; 2000 \cdot 0.00066644\ldots \nonumber\\
  \;&\approx\; 1.33289.
\end{align}

\paragraph{Step 3: total spread.}
\begin{equation}
  \delta^*
  \;=\; 2.5 \;+\; 1.33289
  \;=\; 3.83289\ldots
  \label{eq:ch8_numerical_delta_star}
\end{equation}

\paragraph{Step 4: half-spread.}
\begin{equation}
  \delta^* / 2 \;\approx\; 1.91644.
\end{equation}

\paragraph{Step 5: quotes.}
With $r = 9975$ (long inventory has shifted the reservation price
down by $25$ seashells from the mid $10\,000$),
\begin{align}
  p^{*,\,\mathrm{bid}}
  \;&=\; r - \delta^*/2
  \;=\; 9975 - 1.91644
  \;\approx\; 9973.08, \label{eq:ch8_numerical_bid}\\[0.3em]
  p^{*,\,\mathrm{ask}}
  \;&=\; r + \delta^*/2
  \;=\; 9975 + 1.91644
  \;\approx\; 9976.92. \label{eq:ch8_numerical_ask}
\end{align}

Compare against a symmetric maker centred on the mid $10\,000$ with
the same width $3.833$: that maker quotes $9998.08 / 10\,001.92$.
The A--S quote sits about $25$ seashells lower on both sides, so
the next buy is far more likely to lift our ask at $9976.92$ than a
competitor's at $10\,001.92$; each such fill reduces the long by
one. Meanwhile our bid at $9973.08$ sits well below the mid, so the
next sell is less likely to hit it. That asymmetry is the skew
\cref{sec:ch8_naive} asked for.

\section{The full A--S quoting rule}
\label{sec:ch8_full_rule}

Combine \eqref{eq:ch8_r_closedform} with the optimal-spread
formula from the \nameref{sec:ch8_spread} \texttt{keyfactbox}:
\begin{align}
  p^{*,\,\mathrm{bid}}(t, \Spx, \invq)
  \;&=\; \Spx
         \;-\; \invq\,\riskav\volat^2(\Tmat - t)
         \;-\; \tfrac{1}{2}\delta^*, \label{eq:ch8_bid_final} \\[0.3em]
  p^{*,\,\mathrm{ask}}(t, \Spx, \invq)
  \;&=\; \Spx
         \;-\; \invq\,\riskav\volat^2(\Tmat - t)
         \;+\; \tfrac{1}{2}\delta^*, \label{eq:ch8_ask_final}
\end{align}
with $\delta^*$ from the keyfactbox in \cref{sec:ch8_spread};
\citeCartea{} present the same combined rule in cleaner notation.
Read \eqref{eq:ch8_bid_final}--\eqref{eq:ch8_ask_final} carefully.
The \textbf{centre} $\Spx - \invq\riskav\volat^2(\Tmat - t)$
depends on inventory: long shifts it down, short shifts it up. The
\textbf{width} $\delta^*$ does not depend on inventory at all. The
two pieces are fully separable.

\begin{intuitionbox}
A--S skews your \emph{centre}, never your \emph{width}. The edge per
trade is constant; only the price those trades are centred on moves.
Running long makes both bid and ask cheaper, so the next buy is more
likely to lift the ask (pulling inventory toward zero) and the next
sell less likely to hit the bid (growing it further). Running short
does the opposite. Inventory is controlled by biasing which side
gets hit, not by widening the edge demanded.
\end{intuitionbox}

\section{Parameter intuition and tuning}
\label{sec:ch8_params}

Four parameters, four different sources, four different tuning
stories.

\paragraph{$\volat$: from the data.} Mid volatility is not a
free knob; it is whatever a short-window estimator (e.g.\ an EWMA of
$(\Delta\midp)^2$) reports. Production systems re-estimate $\volat$
at high frequency. On Prosperity estimate it from the previous few
hundred ticks of mid. Units matter: $\volat$ is price per square
root of time, in whichever clock $\Tmat - t$ is measured. Confusing
per-tick with per-session time makes the inventory-risk term land
orders of magnitude off.

\paragraph{$\arrate$: fit from fill data.} The decay $\arrate$
is a property of the book, not a knob. Fit it by logging every quote
at distance $\delta$ and whether it filled before cancellation, then
regressing $\log(\text{fill prob})$ on $\delta$ with slope
$-\arrate$. Liquid crypto markets report $\arrate \in [1, 5]$ per
inverse tick; Prosperity is typically $0.5$--$1.5$ because bot fills
are less price-sensitive than human ones. $\arrate$ is
asset-specific and should be refit when the regime changes.

\paragraph{$\Tmat - t$: given by the session.} On Prosperity a
``session'' is one submitted episode (a few hundred to a few
thousand timesteps). In live equity making it is usually one trading
day; the inventory-risk term shrinks into the close. Trap: as
$\Tmat - t \to 0$, both the inventory term and the reservation-price
skew vanish, so the model loses inventory control in the final
minutes. Production systems either force a flat position or cap
$(\Tmat - t)$ from below.

\paragraph{$\riskav$: the only real knob.} Per
\cref{subsec:ch8_r_worked}, $\riskav$ is the only freely tunable
parameter. Operational rule: pick $\riskav$ so that the inventory
skew at maximum \emph{tolerable} inventory is a meaningful fraction
of $\delta^*$ at that inventory. On Prosperity, with position limit
$50$ and a flat-valued product, start at $\riskav = 10^{-3}$ and
backtest.

\begin{historybox}[title={Historical note: from Ho--Stoll to Avellaneda--Stoikov}]
Inventory-dependent dealer quotes predate A--S. \citet{hostoll1981}
wrote a two-period model where a risk-averse dealer with inventory
$\invq$ quotes asymmetrically to drive expected post-trade inventory
to zero: the same skewing logic, without continuous-time
machinery. \citeGM{} added the informational (adverse-selection)
side four years later, deriving a spread from Bayesian updating
alone. A--S \citeAS{} is the first closed-form unification:
continuous-time, risk-averse, Poisson fills with
distance-dependent intensity. Every subsequent inventory-aware
treatment, including \citet{cartea2015algorithmic}, builds
on the same two-term decomposition.
\end{historybox}

\section{A short reference implementation}
\label{sec:ch8_code}

The whole A--S rule fits in under a dozen lines of Python. The
reference implementation below is a \emph{pure function} from
$(\midp, \invq, \riskav, \volat, \arrate, \Tmat - t)$ to
$(p^{*,\mathrm{bid}}, p^{*,\mathrm{ask}})$, the shape the
Prosperity \texttt{Trader.run(state)} loop wants to call.

\begin{lstlisting}[caption={Avellaneda--Stoikov quoting as a pure function.}]
import math

def as_quotes(mid: float, q: int, gamma: float, sigma: float,
              kappa: float, T_minus_t: float) -> tuple[float, float]:
    """Return (bid, ask) using the Avellaneda-Stoikov formulas.

    mid        : current market mid price
    q          : current inventory (positive=long, negative=short)
    gamma      : risk aversion (1/currency)
    sigma      : volatility of the mid (price / sqrt(time))
    kappa      : exponential decay of fill intensity with distance
    T_minus_t  : remaining session time, same units as sigma
    """
    r = mid - q * gamma * sigma**2 * T_minus_t
    inv_term  = gamma * sigma**2 * T_minus_t
    fill_term = (2.0 / gamma) * math.log(1.0 + gamma / kappa)
    delta_star = inv_term + fill_term
    half = 0.5 * delta_star
    return r - half, r + half
\end{lstlisting}

Two production-use remarks the skeleton does not encode. First,
$\arrate$ should be refit online from the last few hundred
fills; the function assumes the caller passes a fresh value. Second,
$\invq$ should already be clipped to the position limit (e.g.\
$\pm 50$ for RAINFOREST\_RESIN). A well-calibrated bot rarely hits
the clip; a miscalibrated one will produce quotes so skewed that
one side walks off the book. Top Prosperity submissions (see
\cref{sec:ch8_onprosperity}) combine a lightly-skewed base quote
with a hard inventory clip rather than letting A--S do the whole
job.

\section{A--S on Prosperity and on a Goldman desk}
\label{sec:ch8_onprosperity}

\begin{prosperitybox}[title={Prosperity example: RAINFOREST\_RESIN}]
RAINFOREST\_RESIN is the canonical ``static'' Prosperity
product: the fair value of every round is pinned near
$10\,000$ seashells, empirical volatility of the mid over a
single episode sits in $\volat \in [1.5, 2.0]$, and a position
limit of $50$ shares caps directional risk to at most
$\pm 50 \cdot 2 \approx 100$ seashells of mark-to-market swing
over $\pm 1\volat$ moves. Calibrate A--S on a running inventory
$\invq = +4$ with $\arrate = 0.8$, $\riskav = 0.001$,
$\Tmat - t = 500$:
\begin{itemize}
  \item \textbf{Reservation price} at $\volat = 2$:
  $r = 10\,000 - 4 \cdot 0.001 \cdot 4 \cdot 500 = 9992$.
  At $\volat = 1.5$: $r = 10\,000 - 4 \cdot 0.001 \cdot 2.25
  \cdot 500 = 9995.5$.
  \item \textbf{Inventory-risk term} at $\volat = 2$:
  $0.001 \cdot 4 \cdot 500 = 2.0$. At $\volat = 1.5$:
  $0.001 \cdot 2.25 \cdot 500 = 1.125$.
  \item \textbf{Fill-rate term}: $(2/0.001)\ln(1 + 0.001/0.8)
  = 2000 \cdot \ln(1.00125) \approx 2.498$.
  \item \textbf{Total spread}: $\approx 4.498$ at $\volat = 2$,
  $\approx 3.623$ at $\volat = 1.5$.
  \item \textbf{Half-spread}: $\approx 2.249$ or $\approx 1.812$.
  \item \textbf{Quotes at $\volat = 2$}: bid
  $\approx 9989.75$, ask $\approx 9994.25$, both
  \emph{below} the true fair value $10\,000$, skewed down by
  the $+4$ inventory so that the maker is more likely to sell
  than buy.
\end{itemize}
Now compare against the actual top-team policy. The Frankfurt
Hedgehogs' RAINFOREST\_RESIN trader
(\texttt{imc-prosperity-3/FrankfurtHedgehogs\_polished.py},
\texttt{StaticTrader.get\_orders()}) quotes $\texttt{bid\_wall}
+ 1$ and $\texttt{ask\_wall} - 1$: exactly one tick inside
the deepest visible wall prices. The wall mid is always
$10\,000$, so the base-case quotes are the fixed offsets
$9\,999$ bid / $10\,001$ ask (or very near it). Inventory
control is done by a hard clip at the position limit, not by
shifting the centre of the quotes.
\end{prosperitybox}

The comparison is instructive. Pure A--S on RAINFOREST\_RESIN
recommends wider quotes than the top teams use
($\approx 3.6$--$4.5$ vs a fixed $2$-seashell round-trip). A--S is
not wrong; its inventory-risk term is conservatively provisioned for
a product whose true risk is dominated by a known-flat fair value
and a hard position limit, not by mid variance. On RAINFOREST\_RESIN
the mark-to-market variance comes almost entirely from reversals
within a tight range, not accumulating drift, so A--S overstates
the risk and prices it into an overly-wide spread.

The lesson that transfers off Prosperity is the decomposition, not
the literal formula: A--S's reservation-price skew is the right
control knob for inventory management, but when the price process
is not driftless Brownian (flat-value, mean-reverting, news-driven),
the variance term must be recalibrated against the actual source of
risk, or replaced (as top Prosperity teams do) with an external
device like a hard position clip or an EMA fair-value re-anchor.

\begin{gsbox}
A--S (or a close relative) is the backbone of automated internal
market making for single-stock direct-market-access (DMA) flow on
sell-side equity desks. Strats own the calibration pipeline:
$\volat$ refit every five minutes from a rolling realised-vol
window, $\arrate$ refit continuously from live fill rates at each
offset, $\riskav$ set by the risk committee from the desk's daily
VaR (Value at Risk) budget. The trader has one override: a
``desk-breaker'' raising $\riskav$ in real time, tightening
inventory control and widening the fill-rate term at once, pressed
when a name's inventory drifts past an internal soft limit (set
lower than the hard regulatory limit). As a Strats intern, expect
to write the code that reestimates $\arrate$ from overnight fill
logs and to explain why a name's A--S spread widened from $3$ bps
to $5$ bps; the answer is almost always a $\volat$ spike from
overnight news, or an $\arrate$ drop because a flow provider pulled
out and the book thinned. See \cref{ch:strats_role} for the role
itself.
\end{gsbox}

\section{Extensions, limitations, and when A--S breaks}
\label{sec:ch8_breaks}

A--S's closed-form optimality rests on four assumptions, each
violated by some real markets. This section flags the four failure
modes and points at the fixes.

\paragraph{Asymmetric quotes near the position limit.} The
\cref{sec:ch8_spread} derivation assumes $\invq_t$ is unconstrained.
Real makers have hard limits $\bar{\invq}$, and as $|\invq|$
approaches $\bar{\invq}$ the quote that would grow $|\invq|$
further must be suppressed (or moved so far it will not fill).
\citeCartea{} derive the constrained version: near the long limit,
the bid is widened dramatically while the ask stays close to $r$,
so the next fill almost certainly reduces inventory. Practitioners
implement this as a soft wall: multiply the one-sided half-spread
by a factor that grows as $|\invq|/\bar{\invq} \to 1$.

\paragraph{Jump risk and the Brownian assumption.}
\eqref{eq:ch8_sde} makes the mid driftless Brownian. Real mids
jump, cluster (\cref{ch:order_flow} Hawkes), and mean-revert at
short horizons. With drift or jumps, the cancellation of
$\Spx$-linear terms in \cref{subsec:ch8_indifference} fails and the
reservation price picks up a drift term. Fix: replace bare $\Spx$
with $\Spx + \hat{\mu}(\Tmat - t)$ where $\hat{\mu}$ is a
short-horizon drift estimator (EMA of recent returns).

\paragraph{Informed flow.} \eqref{eq:ch8_arrival} makes fill
intensities symmetric in distance to mid, a strong
simplification. With informed traders in the \cref{ch:kyle_gm}
sense, fill intensities are \emph{anti}-correlated with the
direction the mid is about to move. A maker filled on the ask at
$\midp + \delta$ has just received a signal that the mid is more
likely to rise: the adverse-selection effect Glosten--Milgrom
formalised, ignored by A--S. Fix: use an OFI-based estimator
(\cref{ch:order_flow}) and add a drift adjustment to $r$ when OFI
is persistently one-sided.

\paragraph{Queue position.} \Cref{ch:order_flow} established that
queue position is a substantial alpha source: a front-of-queue
resting order at $\bidp$ is worth meaningfully more than a
back-of-queue one at the same price. A--S treats price as the only
decision variable and is blind to queue position. On any
price--time-priority venue with slow quote refresh, the right
operational layer on top of A--S is a queue-aware policy that
avoids cancelling a front-of-queue order unnecessarily, even when
A--S's recommended quote has shifted slightly. \citeBouchaud{}
covers this interaction in detail; queue priority is a separate
concern, applied after the A--S rule.

\begin{pitfallbox}[Don't run vanilla A--S on a trending asset]
A--S's reservation-price skew controls inventory \emph{variance},
not \emph{drift}. On a persistently drifting asset (a trending
stock, a commodity in storage glut, Prosperity's SQUID\_INK during
an Olivia informed-trading window), the machinery will keep
recentering quotes on a mid that is drifting the wrong way, and the
maker accumulates an unhedged position at precisely the wrong time.
Increasing $\riskav$ does not cure it (it widens quotes without
correcting drift); the cure is a separate drift estimator that
shifts the whole quote pair, or refusing to quote when OFI is
persistently one-sided. Vanilla A--S is a pure variance-control
device; plugging it into a drifting environment without a drift
signal is how \cref{sec:ch8_naive}'s pathological run
\eqref{eq:ch8_bad_net} comes back to bite.
\end{pitfallbox}

\section{Key facts to memorise}
\label{sec:ch8_keyfacts}

\begin{itemize}
  \item Naive symmetric quoting earns $\delta$ per trade but is
  exposed to inventory drift $\invq \cdot \Delta \midp$; on
  persistently one-sided flow, drift dominates spread capture.
  \item A--S \textbf{reservation price}
  $r = \Spx - \invq\riskav\volat^2(\Tmat - t)$ is the maker's
  private fair value; quotes are centred on $r$, not on the mid.
  \item A--S \textbf{optimal total spread} is the sum of an
  inventory-risk term $\riskav\volat^2(\Tmat - t)$ and a fill-rate
  term $(2/\riskav)\ln(1 + \riskav/\arrate)$.
  \item As $\riskav \to 0$, the inventory term vanishes and the
  fill-rate term reduces to $2/\arrate$, the independent-sides
  profit-maximising spread.
  \item A--S skews the \emph{centre}, never the \emph{width}.
  Inventory is controlled by biasing which side gets filled.
  \item $\volat$ from data, $\arrate$ from fills, $\Tmat - t$ from
  the session. $\riskav$ is the only dial: tune so
  $\riskav\volat^2(\Tmat - t)$ is comparable to per-trade spread
  capture. Copying from another asset gives a bot that never
  trades.
  \item A--S controls variance, not drift. On trending or
  informed-flow assets, pair with a drift estimator (EMA of
  returns, OFI from \cref{ch:order_flow}).
  \item Top Prosperity teams use fixed offsets from a known flat
  value on RAINFOREST\_RESIN with a hard position clip; A--S is
  the right lens but overstates risk on that product.
  \item A--S is the default automated maker on live sell-side
  equity desks; Strats own $\volat, \arrate$ calibration, trader
  owns the $\riskav$ override.
  \item Ho--Stoll (1981) is the predecessor; Glosten--Milgrom
  (1985) is the parallel adverse-selection development; A--S
  (2008) is the continuous-time unification.
\end{itemize}
